\section{Recursive Bisection Method}

Recursive graph bisection, originally developed for distributing computational
workloads across parallel processors~\cite{karypis1998metis}, can produce a
reproducible geographic benchmark for redistricting. The method represents a
state as a graph, balances population while partitioning the graph, and records
the resulting assignments and parameters for independent review.

\subsection{Benchmark Execution}

The historical national pipeline used census tracts as graph nodes and tract
population as the vertex weight. Edges connect units that share a published
geographic boundary. In the default geographic profile, the edge weight is the
shared-boundary length from the applicable TIGER/Line vintage. METIS minimizes
the weighted cut while recursively allocating the required number of districts.

This profile excludes election results, party registration, incumbency,
candidate identity, race, and community submissions from benchmark
generation. That exclusion is auditable from the configuration and manifest,
but it does not make the profile value-free. Tract resolution, recursive
structure, boundary weighting, population tolerance, random seed, and engine
version can all affect the output.

The governed operational profile is NRS v0.3. It uses 15-digit Census blocks,
versioned State adjacency contexts, a fixed floor/ceiling recursive seat
schedule, manifest-derived seeds, deterministic candidate refinement, and
one-based canonical district labels. For each Census cycle, the published
package binds the executable, input inventory, standard and legal profiles,
State assignment/tree manifests, national proof-coverage summary, and an
independent verification pass. NRS v0.3 is a later research profile; it does
not silently amend the v0.1 schedule incorporated by the candidate model
statute.

Race-conscious VRA-aligned weighting and partisan-weighted analysis are
separate research or case-specific profiles. They are not part of the default
geographic benchmark and must not be used to describe that benchmark as both
race-blind and VRA-compliant.

\subsection{Reproducibility And Configuration}

A reproducible run must identify:

\begin{itemize}
\item census and TIGER/Line releases;
\item geographic-unit index and adjacency definition;
\item district count and population source;
\item structure, edge-weight, search, and engine profiles;
\item balance settings and seed;
\item source commit, compiler, dependency lock, and binary provenance; and
\item canonical hashes linking assignments, analysis, and reports.
\end{itemize}

The current project supports several structures, weight signals, and search
modes. Results from one profile do not establish properties of another. The
candidate national-standard decision record selects standard recursive
bisection, geographic boundary weights, and one precommitted content-derived
seed as a comparison benchmark. Seed wiring and block-level operational
conformance are implemented and verified nationally in NRS v0.3. The earlier
tract outputs analyzed elsewhere in this paper remain separate research
artifacts; they are not NRS v0.3 certificates and cannot supply missing
demographic, partisan, or community evidence for the block-level assignments.
Geometric compactness is supplied only by the later precommitted 2020 bakeoff,
which dissolves NRS and official CD118 assignments over the same retained
Census-block polygons. It measures block-projected geometry rather than the
official plan's original polygon linework.

\subsection{Sequential Public-Cut Ceremony}

For a certified recursive plan, precommitment can be made observable rather
than merely asserted after completion. A genesis artifact binds the rule
profile, root instance, canonical floor/ceiling district tree, review interval,
witness threshold, and challenge policy. The root cut is published after one
review window; subsequent non-leaf cuts are released by tree depth and ordered
by binary path. Each release includes its split instance, certificate, proof,
timestamp, witness receipts, and prior-round hash.

The public verifier checks each disclosed prefix without access to unreleased
descendants, including exact parent-to-child universe derivation. Completion
binds the final certified tree; a sustained challenge instead leaves an
immutable halted transcript. The implementation checks receipt presence and
uniqueness, not an external service's signature or clock, so production claims
require separately validated transparency-log or trusted-timestamp adapters.

Prospective bakeoffs use a separate procedural scoreboard for proof coverage,
prefix verification, ceremony status, transcript size, verifier time, and
negative-test resistance. These fields are not blended with map-quality
metrics into a single winner. Existing bakeoffs remain governed by their
original frozen protocols and carry no retroactive ceremony claim.

\subsection{Population And Contiguity}

Population equality is reported as an observed property of each output, not as
a claim that one configured tolerance is always legally sufficient. For
congressional districts, the controlling legal principle is equality as nearly
exact as practicable. A complete evidence package must report ideal population,
maximum and total deviation, and any post-partition repair.

Contiguity is evaluated against the published adjacency graph. Water, islands,
point contacts, slivers, and bridge corrections can alter that graph and must be
versioned rather than silently repaired.

\subsection{Performance And Evidence Scope}

The Rust implementation makes national experimentation practical, but runtime
depends on state size, engine, data cache, profile, and search mode. Historical
benchmark tables in this paper are engineering measurements tied to their named
environment; they are not universal runtime guarantees.

Ensemble generation is distinct from producing one benchmark plan. A single
run can be reproduced and compared, but it cannot define the distribution of
valid alternatives. Ensemble percentile claims require archived real traces,
multiple chains, convergence diagnostics, immutable election inputs, and
uncertainty reporting.

\subsection{Input-Exclusion Defense}

The defensible procedural claim is narrow: a run whose manifest binds only
population and geographic inputs cannot directly target partisan outcomes
during execution using partisan data that it never receives. This is an
input-exclusion defense, not an impossibility defense.

Partisan consequences may still arise from political geography, and upstream
choices can be manipulated if they are selected after observing outcomes. A
credible standard therefore precommits assignment-affecting profiles, publishes
the benchmark, reports sensitivity, and discloses every baseline-to-final
modification. Human legal and community judgment remains necessary.
